\subsection{What survives a deformation of the Gaussian tensor}

The fourth-root tensor has a useful quadratic extension away
from the CM point.  This also specifies a limitation of a
direct argument by degeneration to the Tate cusp.

\begin{proposition}\label{prop:generic-elliptic-division-norm}
Let
\[
 E_{A,B}:y^2=x^3+Ax^2+Bx,\qquad B(A^2-4B)\ne0,
\]
over a field of characteristic zero.  Thus \((0,0)\) is a
specified nonzero two-torsion point.  For an odd prime
\(\ell\), put \(M=(\ell^2-1)/4\), which is even, and
choose one point in each pair \(\{Q,-Q\}\) of nonzero
\(\ell\)-torsion points.  Define
\[
 F_B(Q)=\frac{x(Q)^2-B}{2y(Q)},\qquad
 D_\ell=B^{-M/2}\prod_{Q\bmod\{\pm1\}}F_B(Q).
\]
Then
\begin{equation}\label{eq:generic-elliptic-norm-square}
                          D_\ell^2=\frac{\chi_{-4}(\ell)}{\ell}.
\end{equation}
Changing representatives changes \(D_\ell\) only by a sign.
It is unchanged by the coordinate scaling
\((x,y,A,B)\mapsto(u^2x,u^3y,u^2A,u^4B)\).

For \(A=0,B=-1\), group the representatives in pairs
\(Q,\iota Q\), one such pair for each order-four CM orbit,
where \(\iota(x,y)=(-x,iy)\).  Then the tensor
\(C_\ell=\prod_{Q\bmod\langle\iota\rangle}F_{-1}(Q)\)
satisfies, with those representative choices,
\begin{equation}\label{eq:quadratic-norm-is-square-of-cm-norm}
                                  D_\ell=C_\ell^2.
\end{equation}
\end{proposition}
\begin{proof}
The addition law gives the duplication identity
\[
                         x(2Q)=\frac{(x(Q)^2-B)^2}{4y(Q)^2}.
\]
The denominator and numerator of \(F_B(Q)\) are nonzero
for nonzero odd-order \(Q\): otherwise \(Q\) is
two-torsion, or \(2Q=(0,0)\), respectively.

The odd division polynomial has the form
\[
                \psi_\ell(x)=\ell
                      \prod_{Q\bmod\{\pm1\}}(x-x(Q)).
\]
We check its constant coefficient.  The initial values at
\((0,0)\) are \(\psi_1=1\), \(\psi_3=-B^2\), and
every even division polynomial vanishes there.  The usual
addition-law recursion
\[
             \psi_{2m+1}=\psi_{m+2}\psi_m^3
                              -\psi_{m-1}\psi_{m+1}^3
\]
then gives, for every odd \(n\),
\[
                         \psi_n(0)=(-1)^{(n-1)/2}B^{(n^2-1)/4}.
\]
For clarity, when \(n=4k+1\), the only surviving term is
\(-\psi_{2k-1}\psi_{2k+1}^3\).  Its sign is positive
and its exponent of \(B\) is
\((k^2-k)+3(k^2+k)=4k^2+2k\).
When \(n=4k+3\), the surviving term is
\(\psi_{2k+3}\psi_{2k+1}^3\).  Its sign is negative
and its exponent is
\((k^2+3k+2)+3(k^2+k)=4k^2+6k+2\).
These are precisely the required exponents.

The number of factors in the displayed product for
\(\psi_\ell\) is \(2M\), which is even.  Evaluation
at zero therefore gives
\[
                    \prod_{Q\bmod\{\pm1\}}x(Q)
                         =\frac{\chi_{-4}(\ell)B^M}{\ell}.
\]
Multiplication by two permutes the nonzero torsion pairs.
Squaring the product defining \(D_\ell\), and applying
the duplication identity, proves
\eqref{eq:generic-elliptic-norm-square}.
The sign assertion follows from \(F_B(-Q)=-F_B(Q)\).
Under the coordinate scaling, each factor of \(F_B\)
is multiplied by \(u\); there are \(2M\) factors.
The normalizing term \(B^{M/2}\) is also multiplied
by \(u^{2M}\), proving invariance.

At the CM point,
\(F_{-1}(\iota Q)=-iF_{-1}(Q)\).  Thus
\[
 D_\ell=(-1)^{-M/2}(-i)^M C_\ell^2=C_\ell^2,
\]
because \(M\) is even.  This proves
\eqref{eq:quadratic-norm-is-square-of-cm-norm}.
\end{proof}

The generic identity controls a square root of
\(\chi_{-4}(\ell)/\ell\).  At the Gaussian point,
the desired tensor is one further square root.
After the value of \(D_\ell\) has been fixed, the two
possible values \(C_\ell\) and \(-C_\ell\) remain.
The involution \(C_\ell\mapsto-C_\ell\) preserves
both \eqref{eq:generic-elliptic-norm-square} and
\eqref{eq:quadratic-norm-is-square-of-cm-norm}.
Consequently a compatibility proved only for the generic
quadratic norm does not establish the fourth-root
compatibility needed for the harmonic residue character.

\begin{proposition}\label{prop:no-cm-preserving-tate-degeneration}
A family of smooth elliptic curves in characteristic zero,
equipped with an origin-preserving automorphism of exact
order four, has constant \(j\)-invariant \(1728\).
It cannot have a degeneration to a multiplicative Tate
cusp while retaining that automorphism on the generic
elliptic curve.  If additional marked points are allowed
to collide in a stable compactification, forgetting
those points and stabilizing still has a smooth
elliptic core of \(j\)-invariant \(1728\).
\end{proposition}
\begin{proof}
Over an algebraic closure of the generic field, put the
curve in short Weierstrass form
\(y^2=x^3+ax+b\).  An automorphism fixing the origin
has the form \((x,y)\mapsto(u^2x,u^3y)\).
Exact order four gives \(u^2=-1\).  Preservation of
the equation then forces \(b=0\); smoothness forces
\(a\ne0\).  The formula
\[
                  j=1728\frac{4a^3}{4a^3+27b^2}
\]
gives \(j=1728\).

For a multiplicative degeneration over a trait, a minimal
semistable equation has nonzero reduction of \(c_4\)
and vanishing discriminant.  Hence
\(j=c_4^3/\Delta\) has a pole.  This is incompatible
with the constant value \(1728\).  Equivalently, in
the compactified moduli of one-pointed elliptic curves,
the boundary has \(j=\infty\), whereas the constant
map has image \(1728\).  The same conclusion after
forgetting and stabilizing additional marks proves the
last assertion.
\end{proof}

Thus a direct Tate-cusp comparison that keeps the
order-four orbit decomposition throughout cannot supply
the desired normalization: that fixed locus has no such
cusp.  If the CM automorphism is forgotten before
deforming, the norm naturally displayed above is
\(D_\ell=C_\ell^2\), and an additional comparison is
needed to recover the remaining sign.  This is an
obstruction to that particular degeneration argument,
not a proof that another full-GT comparison is impossible.
